\section{Related work}
\label{sec:related}

\paragraph{Evidence use in policymaking.} A long line of scholarship asks how, and how much,
research evidence informs policy. \citet{weiss1979many} distinguished instrumental, conceptual,
and symbolic uses of research; more recent work emphasizes the organizational and relational
conditions under which evidence is taken up \citep{oliver2022works,bogenschneider2019revisiting}.
This literature is overwhelmingly qualitative or survey-based and rarely measures, at scale,
whether the evidence a policymaker invokes actually supports the claim being made. We provide such
a measure, turning a normative concern about ``evidence use'' into a quantity that can be estimated
from text.

\paragraph{Text as data and legislative communication.} Press releases are a workhorse corpus for
studying how legislators present themselves and frame policy \citep{grimmer2013representational},
and computational text analysis is now standard in political science
\citep{grimmer2013text,monroe2008fightin}. We build on this tradition but shift the unit of
analysis from words or topics to \emph{claims} and their relationship to an external scientific
record---a comparatively underdeveloped target that requires linking political text to the
research literature.

\paragraph{LLMs as annotators.} A rapidly growing body of work shows that large language models can
match or exceed trained crowd workers on text-annotation tasks including relevance and stance
detection, at a fraction of the cost and with higher consistency
\citep{gilardi2023chatgpt,alizadeh2023open,ziems2024can}. This makes an autonomous coding pipeline
plausible. But naive adoption is dangerous: model labels contain non-random error, and substituting
them for human labels biases downstream estimates and breaks inference even at high accuracy
\citep{egami2023dsl}. Two complementary responses have emerged. First, statistical corrections---%
design-based supervised learning \citep{egami2023dsl} and prediction-powered inference
\citep{angelopoulos2023ppi}---recover valid estimates by combining many surrogate labels with a
small gold set. Second, validation protocols such as the alternative-annotator test
\citep{calderon2025alttest} specify \emph{when} a model may be treated as a replacement annotator.
We use both, and we add a practical ingredient: the gold labels themselves are produced by a
high-capacity LLM judge rather than by humans, with independent NLI models as the surrogate, so the
entire procedure is autonomous.

\paragraph{LLM-as-judge and its pitfalls.} Using one model to evaluate another is now common
\citep{zheng2023judging}, but LLM judges exhibit position, verbosity, and self-preference biases
and can inherit the validity problems of whatever they are calibrated against. The standard
mitigation is to avoid relying on a single judge. Our design follows this prescription by pairing a
generative LLM judge with discriminative NLI classifiers that fail in different ways, and by
reporting cross-method agreement explicitly rather than assuming it.

\paragraph{Scientific claim verification.} Determining whether a document supports a claim is the
subject of scientific fact-checking and claim-verification research, typically framed as
natural-language inference over claim--evidence pairs \citep{williams2018broad,yin2019benchmarking,wadden2020fact}.
\citet{koneru2024llm} study whether LLMs can discern evidence for scientific hypotheses in the
social sciences specifically. Our stance step adopts this framing but embeds it in a measurement
pipeline whose output is a debiased \emph{population} quantity (the party-level support rate), not a
per-pair prediction, which is what shifts the validity burden onto PPI/DSL.

\paragraph{Contribution.} Relative to this work we contribute (i) an end-to-end, fully autonomous
pipeline---claim extraction, evidence retrieval, and support assessment---that needs neither human
RAs nor proprietary APIs; (ii) a concrete demonstration that strict-entailment NLI systematically
under-detects scientific support, so that an LLM judge is not a luxury but is required for an
unbiased estimate; and (iii) substantive findings on evidence use in AI policymaking that this
machinery makes estimable for the first time.
